\documentclass[11pt]{article}
\usepackage{fontspec}
\usepackage{polyglossia}
\setdefaultlanguage{russian}
\setotherlanguage{english}
\setmainfont{DejaVu Serif}
\setsansfont{DejaVu Sans}
\setmonofont{DejaVu Sans Mono}
\usepackage{amsmath,amssymb,amsthm,mathtools}
\usepackage{microtype}
\usepackage{geometry}
\usepackage{enumitem}
\usepackage{booktabs}
\usepackage{xcolor}
\usepackage[hidelinks]{hyperref}
\geometry{margin=1in}
\setlength{\parskip}{0.35em}
\setlength{\parindent}{1.2em}

\newtheorem{theorem}{Теорема}[section]
\newtheorem{proposition}[theorem]{Предложение}
\newtheorem{corollary}[theorem]{Следствие}
\newtheorem{lemma}[theorem]{Лемма}
\theoremstyle{definition}
\newtheorem{definition}[theorem]{Определение}
\newtheorem{example}[theorem]{Пример}
\theoremstyle{remark}
\newtheorem{remark}[theorem]{Замечание}

\newcommand{\Sbx}{\mathfrak S}
\newcommand{\Atom}{\operatorname{Atom}}
\newcommand{\Dec}{\operatorname{Dec}}
\newcommand{\MinNest}{\operatorname{MinNest}}
\newcommand{\SCC}{\operatorname{SCC}}

\title{Размышлизмы о sandbox-атомарности с Commander Sol:\\границы композиции, ранг вложенности и безопасные факторизации носителя}
\author{Alex Malachevsky\\\small ORCID: 0009-0008-6009-3196}
\date{28 августа 2026}

\begin{document}
\maketitle

\begin{abstract}
Об атомарности часто говорят так, будто это внутреннее свойство самого элемента. Для частичной, типизированной и, вообще говоря, некоммутативной системы композиций такой вопрос не имеет определённого смысла, пока не заданы допустимые композиции и класс тривиальных факторов. Мы формализуем \emph{composition sandbox} \(\Sbx=(X,\Omega,U)\), определяем левую, правую и двустороннюю атомарность через входящие свидетели частичных операций и извлекаем индуцированное нетривиальное фактор-отношение. Двусторонний \(U\)-атом в точности является точкой с нулевой входящей степенью в этом отношении. Затем мы отделяем локальную границу атомарности от глобальной границы preorder вложенности. Точный критерий имеет графовый характер: атомы совпадают со всеми nesting-минимальными элементами тогда и только тогда, когда каждая минимальная сильно связная компонента является одноточечной и не содержит внутреннего ребра. В well-founded случае фактор-отношение несёт стандартный ординальный ранг, а атомы в точности являются точками ранга ноль.

Далее мы различаем чистое стирание внешней структуры и настоящее отождествление элементов носителя. Pure erasure сохраняет атомарность, если таблицы частичных операций, типизация и \(U\) не меняются. Обычный quotient частичной алгебры атомарность сохранять не обязан: он может уничтожить атом, слив его result-fiber с составным представителем, и может создать атом, поглотив нетривиальный фактор в quotient-класс тривиальных элементов. При triviality reflection атомарность quotient-класса становится универсальным свойством всего result-fiber. Наконец, после перехода к фактор-отношению стандартное forth/back-условие bounded morphism даёт прозрачный достаточный контракт точного сохранения well-founded factor rank и атомарности. Мы не заявляем новизну частичных алгебр, ординальных рангов, bounded morphisms или transfer homomorphisms по отдельности; вклад состоит в собранной FCOA-картине композиционной границы и строгом разделении локальной атомарности, глобальной вложенности, стирания и геометрии, создаваемой quotient-отождествлением.
\end{abstract}

\section{Мотивация и границы утверждений}

Исходный принцип имеет смысл только в точной форме:
\begin{equation}
\boxed{\text{атомарность относительна к заданному sandbox допустимых композиций}.}
\label{eq:thesis}
\end{equation}
Поэтому слово «простое» в этой работе оставляется классическому арифметическому примеру и не используется как определение общего понятия.

Работа продолжает программу FCOA, начатую в \cite{MalachevskyFCOA2026}, где были разделены значения операций, геометрия definedness, внешняя жёсткость и внутренняя восстанавливаемая память. Здесь изучается другой слой структуры: допустимое разложение и вложенность. Настоящая статья не пересматривает опубликованный checkpoint M0--G1--G2 и не использует недоказанные продолжения основной ветки.

Несколько инструментов ниже являются классическими. Частичные алгебры и их конгруэнции давно изучаются \cite{Gratzer1979,GratzerWenzel1967}; ординальный ранг well-founded relations стандартен в теории множеств \cite{Jech2003}; bounded morphisms являются стандартными конструкциями реляционной и модальной логики \cite{BlackburnDeRijkeVenema2001}; transfer homomorphisms дают механизмы lifting факторизаций в классической теории факторизаций \cite{GeroldingerHalterKoch2006}. Наша задача не состоит в переименовании этих теорий. Нас интересует, что происходит после того, как само factor relation канонически извлекается из типизированного семейства частичных композиций и затем подвергается стиранию или quotient-отождествлению.

\section{Composition sandbox}

\begin{definition}[Sandbox композиции]
\emph{Composition sandbox} — это тройка
\begin{equation}
\Sbx=(X,\Omega,U),
\label{eq:sandbox}
\end{equation}
где \(X\) — носитель, \(\Omega\) — объявленное семейство частичных бинарных операций на соответствующих сортах, а \(U\subseteq X\) — объявленный класс тривиальных элементов или результатов. Типизация является частью сигнатуры. Элемент не становится допустимым фактором лишь потому, что принадлежит носителю: он должен реально допускаться в области определения некоторой операции.
\end{definition}

Для \(\omega\in\Omega\) обозначим её область определения через \(D_\omega\subseteq X\times X\). Ассоциативность, коммутативность, единицы, сокращение и замкнутость не предполагаются.

\begin{definition}[Свидетели разложения]
Для \(x\in X\) положим
\begin{equation}
\Dec_{\Sbx}(x)=\{(\omega,a,b):\omega\in\Omega,\ (a,b)\in D_\omega,\ \omega(a,b)=x\}.
\label{eq:dec}
\end{equation}
Свидетель называется левосторонне нетривиальным, если \(a\notin U\), правосторонне нетривиальным, если \(b\notin U\), и двусторонне нетривиальным, если \(a,b\notin U\).
\end{definition}

Соответствующие множества обозначим \(\Dec_U^L(x)\), \(\Dec_U^R(x)\) и \(\Dec_U^{LR}(x)\).

\begin{definition}[Направленная и двусторонняя атомарность]
Объявленный целевой элемент \(x\notin U\) называется
\begin{enumerate}[label=(\roman*)]
\item левым \(U\)-атомом, если \(\Dec_U^L(x)=\varnothing\);
\item правым \(U\)-атомом, если \(\Dec_U^R(x)=\varnothing\);
\item двусторонним \(U\)-атомом, или просто \(U\)-атомом в этой статье, если
\begin{equation}
\Dec_U^{LR}(x)=\varnothing.
\label{eq:atomdef}
\end{equation}
\end{enumerate}
\end{definition}

В частичном или некоммутативном sandbox эти понятия могут расходиться. Например, если \(U=\{u\}\), а единственная клетка имеет вид \(a\star u=x\), то \(x\) является правым \(U\)-атомом, но не левым.

\begin{definition}[Изолированность и неразложимость]
Элемент называется \emph{изолированным}, если он не встречается ни в одной допустимой клетке операции ни как аргумент, ни как результат. Он называется \emph{неразложимым}, если \(\Dec_{\Sbx}(x)=\varnothing\).
\end{definition}

Всегда
\[
\text{изолированный}\Longrightarrow\text{неразложимый}\Longrightarrow U\text{-атом},
\]
но обратные импликации в общем случае неверны.

\section{Индуцированное нетривиальное factor relation}

\begin{definition}[Фактор-отношение]
На нетривиальной части носителя \(X\setminus U\) определим
\begin{equation}
y\triangleleft x
\label{eq:factorrel}
\end{equation}
тогда и только тогда, когда существует двусторонне нетривиальный свидетель \(\omega(a,b)=x\), в котором \(y=a\) или \(y=b\).
\end{definition}

Отношение \(\triangleleft\) не обязано быть транзитивным или антисимметричным. Его рефлексивно-транзитивное замыкание является preorder; факторизация по взаимной достижимости даёт condensation poset сильно связных компонент.

\begin{proposition}[Локальная графовая характеристика]
Для каждого целевого \(x\notin U\)
\begin{equation}
\boxed{x\text{ является }U\text{-атомом}\iff x\text{ имеет входящую степень }0\text{ в }(X\setminus U,\triangleleft).}
\label{eq:indegree}
\end{equation}
\end{proposition}

\begin{proof}
Двусторонне нетривиальный свидетель существует в точности тогда, когда один из его нетривиальных факторов порождает входящее ребро \(\triangleleft\) в \(x\).
\end{proof}

\begin{theorem}[Монотонность по sandbox]
Пусть \(\Sbx_1=(X,\Omega_1,U)\) и \(\Sbx_2=(X,\Omega_2,U)\) имеют одну и ту же типизацию, причём каждая допустимая клетка \(\Sbx_1\) является допустимой клеткой \(\Sbx_2\) с тем же значением. Тогда
\begin{equation}
\boxed{\Atom(\Sbx_2,U)\subseteq\Atom(\Sbx_1,U).}
\label{eq:monotonicity}
\end{equation}
\end{theorem}

\begin{proof}
Для каждого \(x\)
\[
\Dec^{LR}_{\Sbx_1,U}(x)\subseteq\Dec^{LR}_{\Sbx_2,U}(x).
\]
Поэтому атом более богатого sandbox остаётся атомом после ограничения семейства композиций.
\end{proof}

Иными словами, расширение допустимой композиции может уничтожать атомы, но не создавать их; ограничение композиции может создавать атомы, но не уничтожать уже существующие.

\section{Локальная атомарность и глобальная граница вложенности}

Обозначим через \(\SCC(x)\) сильно связную компоненту элемента \(x\) в factor graph. Элемент назовём \emph{nesting-минимальным}, если его SCC минимальна в condensation poset. Через \(\MinNest(\Sbx,U)\) обозначим объединение всех минимальных SCC.

\begin{theorem}[Атом лежит на глобальной нижней границе]
Для любого sandbox
\begin{equation}
\boxed{\Atom(\Sbx,U)\subseteq\MinNest(\Sbx,U).}
\label{eq:atommin}
\end{equation}
\end{theorem}

\begin{proof}
У атома нет входящих рёбер. Следовательно, у него нет self-loop, он не может лежать в нетривиальной SCC, а его одноточечная SCC не имеет входящих рёбер из других компонент. Значит, эта компонента минимальна.
\end{proof}

Обратное утверждение в циклическом sandbox неверно.

\begin{theorem}[Точный критерий минимальной SCC]
\begin{equation}
\boxed{
\Atom(\Sbx,U)=\MinNest(\Sbx,U)
\iff
\text{каждая минимальная SCC одноточечна и не содержит внутреннего ребра}.}
\label{eq:scccriterion}
\end{equation}
\end{theorem}

\begin{proof}
Одноточечная минимальная SCC без внутреннего ребра не имеет ни внутреннего входящего ребра, ни входящего ребра из другой компоненты, поэтому её единственная точка атомарна. Наоборот, в любой нетривиальной strongly connected component каждая вершина имеет внутреннее входящее ребро. Одноточечная SCC с self-loop также имеет входящее ребро. Такие минимальные компоненты не содержат атомов.
\end{proof}

\begin{remark}
Глобальная ацикличность достаточна, но не необходима для \eqref{eq:scccriterion}. Циклы, расположенные строго выше минимального condensation layer, не мешают совпадению локальной и глобальной границ.
\end{remark}

\begin{example}[Минимальная граница без атомов]
Пусть \(U=\{u\}\) и
\[
a\star a=b,\qquad b\star b=a.
\]
Тогда \(\{a,b\}\) является минимальной SCC, но ни \(a\), ни \(b\) не атомарны.
\end{example}

Именно здесь впервые становится видно, почему фраза «атомы — минимальные элементы» требует ремонта: в общей системе частичных композиций глобальной нижней границей является минимальная SCC, а не обязательно отдельный атом.

\section{Well-founded вложенность и ординальная глубина}

\begin{theorem}[Фактор-ранг]
Пусть \(\triangleleft\) well-founded. Тогда стандартная well-founded recursion задаёт единственный ординальный ранг
\begin{equation}
\rho(x)=\sup\{\rho(y)+1:y\triangleleft x\}.
\label{eq:rank}
\end{equation}
Для каждого целевого элемента
\begin{equation}
\boxed{x\text{ является }U\text{-атомом}\iff\rho(x)=0.}
\label{eq:rankzero}
\end{equation}
Кроме того,
\[
y\triangleleft x\Longrightarrow\rho(y)<\rho(x).
\]
\end{theorem}

\begin{proof}
Существование и единственность \eqref{eq:rank} являются стандартным фактом о well-founded relations \cite{Jech2003}. Ранг ноль эквивалентен отсутствию предшественников, а по Предложению~\ref{eq:indegree} это эквивалентно двусторонней \(U\)-атомарности.
\end{proof}

Finite DAG тем самым является только конечным частным случаем трансфинитной картины. При этом well-foundedness сильнее точного критерия \eqref{eq:scccriterion}: оно запрещает любые ориентированные циклы и бесконечные нисходящие цепи, тогда как совпадение atom/minimal ограничивает только нижний condensation layer.

\section{Тривиальный transport и irreducibility}

Сам факт объявления \(U\) тривиальным классом не делает его группой единиц. Поэтому unit-like поведение должно быть восстановлено из реально объявленных операций.

Для нетривиальных элементов \(y,z\) пишем \(y\to_U z\), если существуют \(u\in U\) и допустимая операция, для которых
\[
\omega(u,y)=z
\qquad\text{или}\qquad
\omega(y,u)=z.
\]
Пусть \(\leadsto_U\) — рефлексивно-транзитивное замыкание, и
\[
y\asymp_U z
\iff y\leadsto_U z\text{ и }z\leadsto_U y.
\]

\begin{definition}[Transport-неразложимость]
Нетривиальный целевой элемент \(x\) называется \emph{\(U\)-transport-irreducible}, если каждый его decomposition witness имеет ровно один фактор в \(U\), а второй фактор \(U\)-ассоциирован с \(x\).
\end{definition}

\begin{definition}[\(U\)-coherence]
Sandbox называется \emph{\(U\)-coherent на целевом сорте}, если
\begin{enumerate}[label=(\roman*)]
\item никакая клетка с двумя факторами из \(U\) не даёт нетривиальный целевой результат;
\item в каждом разложении нетривиального целевого результата с одним фактором из \(U\) второй фактор \(U\)-ассоциирован с результатом.
\end{enumerate}
\end{definition}

\begin{theorem}[Совпадение при coherence]
При \(U\)-coherence
\begin{equation}
\boxed{x\text{ является }U\text{-атомом}\iff x\text{ является }U\text{-transport-irreducible}.}
\label{eq:irreducible}
\end{equation}
\end{theorem}

\begin{proof}
Если \(x\) атомарен, каждое его разложение содержит хотя бы один фактор из \(U\). Первое условие coherence запрещает два тривиальных фактора, второе делает оставшийся cofactor transport-ассоциированным с \(x\). В обратную сторону transport-irreducible элемент по определению не имеет разложения с двумя нетривиальными факторами.
\end{proof}

Следовательно, привычное отождествление «атом = irreducible» допустимо только после явного unit-like контракта.

\section{Erasure и инвариантность terminal value fibers}

\begin{theorem}[Инвариантность при pure erasure]
Пусть erasure забывает только внешние имена, порядок или отношения, но сохраняет носитель, типизацию, \(U\), области определения и значения операций. Тогда изолированность, неразложимость, все направленные классы атомарности, factor relation, SCC и любой well-founded factor rank остаются неизменными.
\end{theorem}

\begin{proof}
Все operation cells и все decomposition witnesses буквально остаются теми же.
\end{proof}

Таким образом, потеря внешнего порядка может резко увеличить группу автоморфизмов и при этом вообще не изменить композиционную границу.

\begin{theorem}[Инвариантность terminal value fibers]
Пусть два sandbox имеют один и тот же носитель, \(U\), области определения операций и один и тот же набор клеток, результаты которых лежат в выбранном atomicity-target sort. Они могут различаться только разбиением terminal-result cells по анонимным terminal outputs. Тогда их классы атомов на target sort совпадают.
\end{theorem}

\begin{proof}
Атомарность \(x\) зависит только от двусторонне нетривиальных клеток с результатом \(x\). Переразбиение terminal-output fibers ни одной такой клетки не меняет.
\end{proof}

Это отделяет atomicity активного результата от value-fiber rigidity, изучавшейся в предыдущей линии FCOA \cite{MalachevskyFCOA2026}.

\section{Обычный quotient может менять атомарность}

Пусть
\[
q:X\twoheadrightarrow\bar X
\]
— sort-respecting congruence quotient, совместимый с частичными операциями и использующий existential representative semantics для quotient cells, а \(\bar U=q(U)\).

В этом разделе quotient partial operation понимается в следующем явном смысле. Для quotient-классов \(\bar a,\bar b,\bar z\)
\begin{equation}
\boxed{
\bar\omega(\bar a,\bar b)=\bar z
\iff
\exists a'\in q^{-1}(\bar a)\;\exists b'\in q^{-1}(\bar b)\;\exists z'\in q^{-1}(\bar z)
\text{ такие, что }\omega(a',b')=z'.}
\label{eq:quotsem}
\end{equation}
Условие congruence/compatibility гарантирует, что если такая quotient-cell определена, то результирующий класс \(q(z')\) не зависит от выбора представителей, дающих свидетель. Мы не требуем, чтобы каждая пара представителей классов \((\bar a,\bar b)\) лежала в исходной области определения: используемая семантика намеренно existential.

\begin{theorem}[Точный quotient witness criterion]
Класс \(q(x)\) неатомарен тогда и только тогда, когда существуют \(a,b,z\in X\) и \(\omega\in\Omega\), такие что
\begin{equation}
q(z)=q(x),\qquad \omega(a,b)=z,
\qquad q(a),q(b)\notin\bar U.
\label{eq:quotwitness}
\end{equation}
\end{theorem}

\begin{proof}
Двусторонне нетривиальный quotient witness по определению quotient semantics представлен допустимой исходной клеткой, оба фактор-класса которой нетривиальны, а результат принадлежит нужному quotient result fiber. Это и есть \eqref{eq:quotwitness}.
\end{proof}

\begin{example}[Загрязнение result-fiber]
Пусть
\[
X=\{u,a,b,x,y\},\qquad U=\{u\},
\]
и единственная нетривиальная клетка равна
\[
a\star b=y.
\]
Элемент \(x\) атомарен, а \(y\) составен. Отождествим \(x\sim y\). Тогда quotient-класс \([x]=[y]\) составен, поскольку свидетель для \(y\) спускается в этот класс. Атомарный представитель потерял атомарность только потому, что в его result-fiber попал составной представитель.
\end{example}

\begin{example}[Triviality collapse создаёт атом]
Пусть
\[
X=\{u,a,b,x\},\qquad U=\{u\},
\]
и
\[
a\star b=x,\qquad u\star b=x.
\]
Элемент \(x\) не атомарен из-за первой клетки. Если отождествить \(a\sim u\), то \([a]=[u]\in\bar U\). После quotient у \([x]\) остаётся только представление с тривиальным фактором, и \([x]\) становится атомом. Следовательно, quotient может создавать атомы, поглощая нетривиальный фактор в тривиальный класс.
\end{example}

\begin{definition}[Triviality reflection]
Quotient называется \emph{triviality-reflecting}, если
\begin{equation}
\boxed{q^{-1}(\bar U)=U.}
\label{eq:trivreflect}
\end{equation}
\end{definition}

\begin{theorem}[Универсальный fiber-критерий]
При triviality reflection
\begin{equation}
\boxed{
q(x)\in\Atom(\bar\Sbx,\bar U)
\iff
q^{-1}(q(x))\subseteq\Atom(\Sbx,U).}
\label{eq:fibercriterion}
\end{equation}
\end{theorem}

\begin{proof}
Из \eqref{eq:trivreflect} следует \(q(a)\notin\bar U\iff a\notin U\). Поэтому двусторонне нетривиальный quotient witness существует в точности тогда, когда некоторый представитель результата в fiber класса \(q(x)\) имеет исходный двусторонне нетривиальный witness. Отрицание этого утверждения даёт \eqref{eq:fibercriterion}.
\end{proof}

Это один из главных результатов quotient-части: при triviality reflection atomhood является не свойством выбранного представителя, а универсальным свойством всего result-fiber.

\section{Безопасные quotient maps на factor frame}

Обычная quotient compatibility вместе с triviality reflection даёт forward property:
\begin{equation}
y\triangleleft x\Longrightarrow q(y)\,\bar\triangleleft\,q(x).
\label{eq:forth}
\end{equation}
Недостающим условием является стандартный back clause.

\begin{definition}[Bounded morphism factor frame]
Пусть triviality reflection выполнено. Индуцированное отображение
\[
q:(X\setminus U,\triangleleft)\twoheadrightarrow
(\bar X\setminus\bar U,\bar\triangleleft)
\]
называется \emph{surjective bounded morphism factor frames}, если оно удовлетворяет \eqref{eq:forth} и
\begin{equation}
\boxed{
\bar y\,\bar\triangleleft\,q(x)
\Longrightarrow
\exists y\in q^{-1}(\bar y)\text{ такое, что }y\triangleleft x
}
\label{eq:back}
\end{equation}
для каждого нетривиального представителя \(x\).
\end{definition}

Квантор по каждому представителю result-class принципиален. Условие \eqref{eq:back} является в точности стандартным back condition bounded/p-morphism после того, как factor relation рассматривается как accessibility relation \cite{BlackburnDeRijkeVenema2001}. Во внутренних заметках ветки это условие называлось coherent predecessor lifting, но отдельная новизна общего morphism concept здесь не заявляется.

\begin{theorem}[Сохранение well-foundedness]
Если исходное \(\triangleleft\) well-founded, а индуцированное отображение factor frames является surjective bounded morphism, то \(\bar\triangleleft\) также well-founded.
\end{theorem}

\begin{proof}
Пусть в quotient существует бесконечная нисходящая цепь
\[
\bar x_0\succ\bar x_1\succ\bar x_2\succ\cdots.
\]
Выберем любой представитель \(x_0\) класса \(\bar x_0\). Последовательное применение back clause рекурсивно поднимает цепь к
\[
x_0\succ x_1\succ x_2\succ\cdots
\]
в исходном factor relation, что противоречит well-foundedness.
\end{proof}

\begin{theorem}[Точное сохранение ранга]
При тех же условиях, если \(\rho\) и \(\bar\rho\) — well-founded factor ranks исходной структуры и quotient, то
\begin{equation}
\boxed{\bar\rho(q(x))=\rho(x)}
\label{eq:rankpreserve}
\end{equation}
для каждого \(x\notin U\).
\end{theorem}

\begin{proof}
Используем well-founded induction по \(x\). Предположим, что
\[
\bar\rho(q(y))=\rho(y)
\]
для каждого предшественника \(y\triangleleft x\). По определению quotient rank
\begin{equation}
\bar\rho(q(x))
=
\sup_{\bar y\,\bar\triangleleft\,q(x)}
\bigl(\bar\rho(\bar y)+1\bigr).
\label{eq:quotranksup}
\end{equation}

Сначала получим верхнюю оценку. Возьмём произвольный \(\bar y\bar\triangleleft q(x)\). По back condition существует \(y\triangleleft x\) такое, что \(q(y)=\bar y\). По индуктивной гипотезе
\[
\bar\rho(\bar y)=\bar\rho(q(y))=\rho(y),
\]
поэтому
\[
\bar\rho(\bar y)+1=\rho(y)+1\le\rho(x).
\]
Беря supremum в \eqref{eq:quotranksup}, получаем
\begin{equation}
\bar\rho(q(x))\le\rho(x).
\label{eq:rankupper}
\end{equation}

Для нижней оценки используем определение исходного ранга:
\[
\rho(x)=\sup_{y\triangleleft x}\bigl(\rho(y)+1\bigr).
\]
Для каждого \(y\triangleleft x\) forth condition даёт \(q(y)\bar\triangleleft q(x)\), а индуктивная гипотеза — \(\rho(y)=\bar\rho(q(y))\). Следовательно,
\[
\rho(y)+1
=
\bar\rho(q(y))+1
\le
\bar\rho(q(x)).
\]
Беря supremum по всем исходным предшественникам, получаем
\begin{equation}
\rho(x)\le\bar\rho(q(x)).
\label{eq:ranklower}
\end{equation}
Совместно \eqref{eq:rankupper} и \eqref{eq:ranklower} дают \eqref{eq:rankpreserve}.
\end{proof}

\begin{corollary}[Точное сохранение атомарности]
При условиях предыдущей теоремы
\begin{equation}
\boxed{x\text{ является }U\text{-атомом}\iff q(x)\text{ является }\bar U\text{-атомом}.}
\end{equation}
\end{corollary}

\begin{proof}
В well-founded factor relation атомарность эквивалентна рангу ноль; остаётся применить \eqref{eq:rankpreserve}.
\end{proof}

Bounded morphism здесь является достаточным safety contract для точного сохранения ранга. Мы не утверждаем, что он необходим для каждого quotient, который случайно сохраняет атомарность или well-foundedness.

\section{Связь с transfer homomorphisms}

Transfer homomorphisms в теории факторизаций моноидов поднимают существенно больше структуры, чем одношаговое predecessor incidence: они предназначены для переноса факторизационной структуры и включают lifting factorization of the image \cite{GeroldingerHalterKoch2006}. Поэтому factor-frame theorem настоящей статьи не является конкурентной «общей теорией факторизаций».

Связь концептуальна. Из произвольных типизированных частичных операций мы сначала извлекаем только factor relation. На этом реляционном уровне back clause bounded morphism является ровно тем условием, которое поднимает предшественников. В более богатом ассоциативном monoid setting transfer homomorphism сохраняет соответственно более богатые данные о факторизациях. Такое сопоставление помещает FCOA-конструкцию внутрь известной математики и предотвращает завышение novelty claim.

\section{Классические простые как один sandbox}

Рассмотрим
\begin{equation}
\Sbx_{\mathbb N}=(\mathbb Z_{>0},\{\cdot\},\{1\}).
\label{eq:classical}
\end{equation}
Для \(n>1\) двусторонне нетривиальное разложение — это в точности
\[
ab=n,\qquad a>1,\quad b>1.
\]
Поэтому двусторонние \(\{1\}\)-атомы совпадают с обычными простыми числами.

Однако классический sandbox обладает дополнительной структурой. Единица \(1\) является настоящей единицей, собственные нетривиальные множители меньше результата, а factor relation well-founded относительно обычной величины числа. Поэтому несколько понятий, расходящихся в общем sandbox, здесь совпадают. Уникальность факторизации остаётся существенно более сильной теоремой и из определения атома не следует.

Правильное направление интерпретации таково:
\[
\boxed{\text{классические простые — один частный sandbox атомарности, а не её общее определение}.}
\]

\section{Чему учат конечные контрпримеры}

Конечные witnesses ветки разделяют логически разные эффекты: один и тот же carrier при неизменном \(U\) может сделать точку атомарной или составной только за счёт изменения admissible cells; indecomposable не обязано означать isolated; atomic не обязано означать indecomposable; left/right atomicity могут расходиться; минимальная SCC может не содержать атомов; изменение только terminal value fibers способно менять rigidity, не меняя active atomicity; pure erasure сохраняет атомарность, тогда как quotient identification — нет; ordinary quotient имеет два независимых механизма повреждения атомарности; bounded factor morphism устраняет эти патологии в well-founded rank setting.

Эти примеры специально малы. Их задача — не численный эксперимент, а логическое разделение понятий, которые в привычной арифметике часто автоматически совпадают.

\section{Ограничения и дисциплина утверждений}

Результаты статьи относятся к типизированным семействам частичных бинарных композиций и factor relation, индуцированному двусторонне нетривиальными decomposition witnesses. Они не доказывают существование atomic decomposition для каждого элемента, уникальность факторизации, канонический выбор \(U\), ассоциативность или коммутативность, не объявленные в сигнатуре, сохранение атомарности произвольными quotient-identifications или необходимость bounded-morphism condition для всех rank-preserving quotients.

Также не заявляется новизна partial algebras, strong congruences, well-founded ordinal rank, bounded morphisms или transfer homomorphisms как общих теорий. Специфический результат этой ветки FCOA состоит в собранной boundary-картине:
\begin{equation}
\boxed{
\begin{aligned}
\text{локальная граница композиции} &\;=\; \text{атомарность},\\
\text{глобальная граница nesting} &\;=\; \text{минимальный SCC-layer},\\
\text{well-founded глубина} &\;=\; \text{ординальный factor rank},\\
\text{обычное отождествление} &\;\text{ может создавать factor geometry},\\
\text{bounded factor morphism} &\;\text{ сохраняет rank-layer точно}.
\end{aligned}}
\label{eq:synthesis}
\end{equation}

\section{Заключение}

Фраза «атомарность есть граничное состояние композиции» становится математической только после задания sandbox. \(U\)-атом — это локальная точка с нулевым входящим factor incidence. Такая локальная граница всегда лежит в минимальном condensation layer, но совпадает со всей глобальной границей только тогда, когда каждая минимальная SCC является одноточечной и не содержит внутреннего ребра. При well-founded nesting то же утверждение приобретает ранговую форму: атомы — точки ранга ноль.

Различие erasure и quotient столь же принципиально. Pure erasure вообще не меняет decomposition witnesses. Quotient identification может менять их косвенно, смешивая result fibers или меняя принадлежность фактор-классов к \(U\). При triviality reflection atomhood становится универсальным свойством всего result-fiber. Если индуцированное отображение factor frames дополнительно удовлетворяет стандартному back clause bounded morphism, well-founded factor rank сохраняется точно.

Полученная иерархия даёт язык для исследования границ композиции без импорта обычной делимости и без превращения привычных арифметических совпадений в аксиомы общей теории.

\section*{Примечание о термине Commander Sol}
Commander Sol — обозначение AI-assisted исследовательской роли в данной программе. Математические утверждения оформлены как определения, теоремы и доказательства и подчинены той же дисциплине claims, что и остальные работы FCOA.

\bibliographystyle{plain}
\bibliography{references}

\end{document}
